\section{本章小结}
\label{sec:dfam-summary}

章首的问题是：在挑选一个分布之前，怎样说清允许哪些分布、怎样给它们编号。第一节把这两件事分开。候选范围由支撑、矩约束和相对参考测度的绝对连续性规定，它们限制的是概率规律本身；三面骰子的候选集合是三角形中的线段，给定均值的密度却多到不能用有限个系数线性表示。参数化是从参数到分布的映射，可辨识性问的是这个映射是否为单射。重复参数、只以组合出现的参数和标签置换都会破坏可辨识性；可辨识只是分布层面的唯一性，不等于有限样本能准确分辨参数。

后四节给出四种构造。它们固定的对象与可调的参数不同，表~\ref{tab:dfam-constructions}把它们并排列出。

\begin{BookInlineTable}
  \centering
  \caption{四种构造分布族的方式}
  \label{tab:dfam-constructions}
  \begin{tabular}{@{}llll@{}}
    \toprule
    构造 & 固定的对象 & 随参数变化的部分 & 需要核对的条件\\
    \midrule
    推前 & 基础分布$P_0$ & 映射$T_{\bm{\theta}}$ & 密度是否存在\\
    混合 & 核或成分族 & 混合测度、成分参数 & 表示是否唯一\\
    Copula & 依赖结构$C$ & 边缘分布$F_i$ & 边缘是否连续\\
    指数族 & $h\nu$与$\bm{T}$ & 自然参数$\bm{\theta}$ & 参数域与最小性\\
    \bottomrule
  \end{tabular}
\end{BookInlineTable}

推前族把随机性放在固定的基础分布里，参数只控制映射。位置尺度族中分布函数与分位数按仿射规律变化，矩在存在时随之变化，标准化后的形状保持不变；固定基础分布后，非退化的基础变量使位置与尺度可辨识。非线性可逆映射改变形状，低维随机性进入高维空间使支撑退化而失去密度，离散化产生原子。只有同维可逆光滑的映射才能直接用Jacobian公式写出密度，一般的推前族要用推前测度和基础空间上的期望来描述。

混合族按随机选出的成分产生结果，有限混合是概率的加权平均，而不是取值的加权平均。两个高斯成分足够分开时出现双峰，这是成分族自身没有的形状；改权重只在固定形状之间重新分配质量，改成分参数才改变形状。连续混合把求和换成对混合测度的积分：正态尺度混合的峰度不小于$3$，多元时共同尺度使分量不相关却不独立；Poisson–Gamma混合得到方差更大的负二项分布。混合表示有标签置换、重复成分与零权重三种冗余，成分重叠本身并不是冗余：已知方差的高斯位置混合除标签置换外可辨识，Bernoulli成分的混合却完全不可辨识。

Copula把联合分布拆成边缘与依赖。Sklar定理说任何Copula可以与任何边缘拼成联合分布，任何联合分布都有Copula；边缘连续时，概率积分变换把各分量送到均匀刻度，Copula唯一，并且在分量的严格递增变换下不变。边缘有原子时，Copula只在边缘分布函数值域的乘积上确定，两枚独立硬币可以配上看起来很不一样的Copula，分布变换仍能保证存在性。

指数族对若干统计量施加指数权重，基准测度与统计量固定，自然参数在使配分积分有限的凸集上变化。参数域内部，对数配分函数的梯度与Hessian给出统计量的均值与协方差，凸性由协方差的半正定性解释，严格凸性要求统计量之间没有仿射关系。最小表示下参数可辨识；正则族中均值映射是从自然参数域到一个开集的连续可微双射，但这个开集不含可行均值的边界，参数趋于无穷时分布可能趋于族外的点质量，非正则族的边界参数也可能没有有限均值。

这四种构造可以描述同一个分布族。高斯族既是标准高斯的位置尺度族，又是以$(x,x^{2})$为统计量的指数族；多元高斯分布既是标准高斯向量的仿射推前，也可以拆成高斯边缘加上高斯Copula；Laplace分布既是位置尺度族，又是正态尺度混合。选用哪一种表示，取决于要回答的问题：生成样本时推前最方便，解释异质来源时混合最自然，分开处理边缘与依赖时用Copula，需要用矩刻画分布时用指数族。

本章只建立了候选族的表示与可辨识性。给定约束后在全体允许的分布中选出哪一个，指数族会作为最优解的形式出现，这是第~\ref{chap:distribution-selection}章的内容；把候选限制在某个族之后，怎样在族内逼近一个给定的分布、误差有多大，由第~\ref{chap:distribution-approximation}章研究。
